\subsection{Das Pascalsche Dreieck}
\begin{LA}{Pascalsches Dreieck}{pascalschesdreieck}
    Zur Vereinfachung von Termen der Form \((a+b)^n\) kann man das sogenannte \textit{Pascalsche Dreieck} verwenden. In dieser Lernaufgabe werden Sie dieses Dreieck selbst entdecken und als Hilfsmittel zur Bestimmung höherer Potenzen anwenden.

    \begin{enumerate}
        \item Berechnen Sie die folgenden Potenzen:
              \begin{enumerate}
                  \item \((a+b)^1\)
                  \item \((a+b)^2\)
                  \item \((a+b)^3\)
                  \item \((a+b)^4\)
              \end{enumerate}

        \item Ergänzen Sie mit den Resultaten aus Aufgabe 1) das folgende Dreieck, indem Sie wie in den ersten beiden Zeilen vorgemacht, die fehlenden Koeffizienten einfügen.
              \[
                  \begin{aligned}
                      (a+b)^1 & = \phantom{a^4+4a^3+6a^2}{\color{red}1}\,a+{\color{red}1}\,b                              \\
                      (a+b)^2 & = \phantom{a^4+4a^3+}{\color{red}1}\,a^2+{\color{red}2}\,ab+{\color{red}1}\,b^2           \\
                      (a+b)^3 & = \phantom{a^4+4}a^3+\phantom{1}\quad a^2b+\phantom{1}\quad ab^2+\quad b^3                \\
                      (a+b)^4 & = \quad a^4+\phantom{1}\quad a^3b+\phantom{1}\quad a^2b^2+\phantom{1}\quad ab^3+\quad b^4
                  \end{aligned}
              \]

        \item Untersuchen Sie das Resultat aus Aufgabe 2 und erklären Sie, wie sich die Koeffizienten einer Zeile aus jenen der Zeile davor berechnen lassen. Wie müssten demnach die nächsten drei Zeilen lauten?
              \begin{enumerate}
                  \item \((a+b)^5\)
                  \item \((a+b)^6\)
                  \item \((a+b)^7\)
              \end{enumerate}

              Lassen Sie das Resultat von der Lehrperson kontrollieren.

        \item Berechnen Sie die folgenden Potenzen:
              \begin{enumerate}
                  \item \((a-b)^1\)
                  \item \((a-b)^2\)
                  \item \((a-b)^3\)
                  \item \((a-b)^4\)
              \end{enumerate}

        \item Ergänzen Sie mit den Resultaten aus Aufgabe 4) das folgende Dreieck, indem Sie wie in der ersten Zeile vorgemacht, die fehlenden Koeffizienten einfügen.
              \[
                  \begin{aligned}
                      (a-b)^1 & = \phantom{a^4+4a^3+6a^2}{\color{red}1}\,a-{\color{red}1}\,b                              \\
                      (a-b)^2 & = \phantom{a^4+4a^3+}{\color{red}1}\,a^2-{\color{red}2}\,ab+{\color{red}1}\,b^2           \\
                      (a-b)^3 & = \phantom{a^4+4}a^3-\phantom{1}\quad a^2b+\phantom{1}\quad ab^2-\quad b^3                \\
                      (a-b)^4 & = \quad a^4-\phantom{1}\quad a^3b+\phantom{1}\quad a^2b^2-\phantom{1}\quad ab^3+\quad b^4
                  \end{aligned}
              \]

        \item Untersuchen Sie, wie sich die Koeffizienten im Dreieck entwickeln, und finden Sie eine Regelmässigkeit. Bestimmen Sie anschliessend \textbf{nur mit Hilfe des Dreiecks}
              \begin{enumerate}
                  \item \((a-b)^5\)
                  \item \((a-b)^6\)
                  \item \((a-b)^7\)
              \end{enumerate}

              Lassen Sie das Resultat von der Lehrperson kontrollieren.
    \end{enumerate}
\end{LA}

\clearpage
Zur Vereinfachung von Termen der Form \((a+b)^n\) kann man das sogenannte
\textit{Pascalsche Dreieck} verwenden. Dieses Zahlendreieck ist nach dem
französischen Mathematiker Blaise Pascal (1623--1662) benannt\footnote{
    Dieses Zahlendreieck war aber schon früher bekannt. Entsprechend hat das
    Dreieck in anderen Ländern auch andere Namen -- jeder bezieht sich auf einen
    Mathematiker aus der entsprechenden Gegend! In Italien wird es
    Tartaglia-Dreieck genannt, in China nennt man es das Yang-Hui-Dreieck und im
    Iran spricht man vom Chayyam-Dreieck.}.

\begin{bsp}{Anwendung des Pascalschen Dreiecks}{pascalbeispiel1}
    Zur Vereinfachung von \((a+b)^6\) verwenden wir die sechste Zeile des
    Pascalschen Dreiecks:
    \vspace*{10 cm}
    % \[
    %     \begin{array}{ccccccccccccc}
    %           &   &   &   &    &    & 1                            \\
    %           &   &   &   &    & 1  &    & 1                       \\
    %           &   &   &   & 1  &    & 2  &    & 1                  \\
    %           &   &   & 1 &    & 3  &    & 3  &    & 1             \\
    %           &   & 1 &   & 4  &    & 6  &    & 4  &   & 1         \\
    %           & 1 &   & 5 &    & 10 &    & 10 &    & 5 &   & 1     \\
    %         1 &   & 6 &   & 15 &    & 20 &    & 15 &   & 6 &   & 1
    %     \end{array}
    % \]

    % Die Koeffizienten von \((a+b)^6\) sind also
    % \[
    %     1,\;6,\;15,\;20,\;15,\;6,\;1.
    % \]
    % Damit erhält man
    % \[
    %     (a+b)^6
    %     =
    %     a^6+6a^5b+15a^4b^2+20a^3b^3+15a^2b^4+6ab^5+b^6.
    % \]
\end{bsp}

\begin{bsp}{Ein zweites Beispiel}{pascalbeispiel2}

    \[
        (x^2-2y)^4
    \]

    \vspace*{5 cm}
    % Wie verwenden die Koeffizienten
    % \[
    %     1,\;4,\;6,\;4,\;1.
    % \]
    % Somit gilt
    % \[
    %     \begin{aligned}
    %         (x^2-2y)^4
    %          & =(x^2)^4
    %         +4(x^2)^3(-2y)
    %         +6(x^2)^2(-2y)^2
    %         +4x^2(-2y)^3
    %         +(-2y)^4                               \\
    %          & =x^8-8x^6y+24x^4y^2-32x^2y^3+16y^4.
    %     \end{aligned}
    % \]
\end{bsp}

\clearpage
\begin{ueb}
    {Pascalsches Dreieck}{pascalschesdreieck}
    \begin{enumerate}
        \item
              \begin{enumerate}
                  \item \((2a+b)^3\)
                  \item \((a-2b)^3\)
                  \item \((3c+1)^4\)
                  \item \((4-d)^4\)
              \end{enumerate}

        \item
              \begin{enumerate}
                  \item \((2p+q)^5\)
                  \item \((-r+10)^5\)
                  \item \(\left(s+\dfrac{1}{2}\right)^4\)
                  \item \(\left(t+\dfrac{1}{t}\right)^4\)
              \end{enumerate}

        \item
              \begin{enumerate}
                  \item \((2x+y)^3-(2x-y)^3\)
                  \item \(\left(z+\dfrac{1}{2}\right)^4-\left(z-\dfrac{1}{2}\right)^4\)
                  \item \(\left(z+\dfrac{1}{2}\right)^4+\left(z-\dfrac{1}{2}\right)^4\)
                  \item \(\left(u-\dfrac{1}{u}\right)^3-\left(u+\dfrac{1}{u}\right)^3\)
              \end{enumerate}

        \item
              \begin{enumerate}
                  \item \((m^2+m)^3\)
                  \item \((n^2+2n)^3\)
                  \item \((p^3-1)^3\)
                  \item \((1-q^4)^4\)
              \end{enumerate}

        \item \textbf{Challenge}
              \begin{enumerate}
                  \item Berechnen Sie mit dem Taschenrechner die Potenzen
                        \[
                            11^0,\;11^1,\;11^2,\;11^3,\;11^4
                            \quad\text{und}\quad 11^5.
                        \]
                        Erklären Sie das Resultat.

                  \item Berechnen Sie zeilenweise die Summe der Zahlen im
                        Pascalschen Dreieck. Erklären Sie das Resultat.
              \end{enumerate}
    \end{enumerate}
\end{ueb}

\clearpage
\emph{Lösungen}
\begin{enumerate}
    \item
          \begin{enumerate}
              \item \(8a^3+12a^2b+6ab^2+b^3\)
              \item \(a^3-6a^2b+12ab^2-8b^3\)
              \item \(81c^4+108c^3+54c^2+12c+1\)
              \item \(d^4-16d^3+96d^2-256d+256\)
          \end{enumerate}

    \item
          \begin{enumerate}
              \item \(32p^5+80p^4q+80p^3q^2+40p^2q^3+10pq^4+q^5\)
              \item \(-r^5+50r^4-1000r^3+10000r^2-50000r+100000\)
              \item \(s^4+2s^3+\dfrac{3}{2}s^2+\dfrac{1}{2}s+\dfrac{1}{16}\)
              \item \(t^4+4t^2+6+\dfrac{4}{t^2}+\dfrac{1}{t^4}\)
          \end{enumerate}

    \item
          \begin{enumerate}
              \item \(24x^2y+2y^3\)
              \item \(4z^3+z\)
              \item \(2z^4+3z^2+\dfrac{1}{8}\)
              \item \(-6u-\dfrac{2}{u^3}\)
          \end{enumerate}

    \item
          \begin{enumerate}
              \item \(m^6+3m^5+3m^4+m^3\)
              \item \(n^6+6n^5+12n^4+8n^3\)
              \item \(p^9-3p^6+3p^3-1\)
              \item \(q^{16}-4q^{12}+6q^8-4q^4+1\)
          \end{enumerate}
\end{enumerate}































